\begin{abstract}

Planning a university degree is a constraint-satisfaction problem that students usually attempt
with a printed curriculum sheet and a spreadsheet. Neither can say which courses are currently
eligible, what a dropped prerequisite invalidates, which elective group should claim a repeated
course, or whether the final-year thesis track is still available. This thesis presents the IU Smart
Study Planner, a web application for students at International University -- VNU HCMC, and reports
what it does and does not establish.

The artifact contributes a program-aware curriculum model in which courses are global while
membership, placement and prerequisite edges are per program, so shared courses are not duplicated
and two programs are never silently merged. Accounts use cookie-based sessions whose role and
program context are re-read on every request, and each mutation carries an expected scope that is
compared inside the writing transaction, turning a stale client into a conflict rather than a silent
error. Completion is server-authoritative, with mandatory prerequisite validation and a
transactional transitive cascade on uncompletion. Grade attempts are immutable and append-only,
feeding a credit-weighted average computed from the highest score per course in decimal arithmetic.
Completion-gated difficulty ratings are combined with a program prior through Bayesian shrinkage and
always displayed with their vote count. A saved-semester editor validates ownership, revision,
membership, prerequisites and grade eligibility in one serializable transaction with immutable
receipts, recoverable drafts and per-course blocker explanations. An administrative surface
simulates resource settings, cohort demand, capacity and two allocation models behind frozen,
replayable contracts executed by an explicit single-job command.

Verification uses fail-first tests against a real PostgreSQL database, one reproducible
continuous-integration gate, local profiling of the allocation models and a browser accessibility
audit. The artifact makes no claim of institutional validation, calendar validity or improved
student outcomes, and reports those limits explicitly rather than filling them.

\end{abstract}
